\section{Sampling from \ComputationActivationNetworks{}}\label{sec:sampling}

In general quantum rejection sampling consist in two steps:
\begin{itemize}
    \item Find an ancilla augmentation of a distribution, such that there is a circuit preparing the q-sample of the augmented distribution
    \item Amplify the amplitude of the ancilla qubits in state $1$ by repeated Grover operators
\end{itemize}

We investigate, how the above circuit encoding schemes can be applied in the preparation of q-samples to \ComputationActivationNetworks{}.
In particular, we build compositions of \computationCircuits{} and \activationCircuits{}, for which conditional distributions coincide with \ComputationActivationNetworks{}.

% Graph-controlled circuits
Both are defined using controlled single qubit gates (see Sections 4.2-3 in \cite{nielsen_quantum_2010}) with ancilla qubits being the target qubits.

\subsection{Preparing ancilla augmented distributions for \ComputationActivationNetworks{}}

The computation network consists already of directed cores and therefore is already a Bayesian network.
The activation network however needs ancilla augmentation.

\begin{figure}[t]
    \begin{center}
        \input{tikz_pics/sampling/can_elementary_augmentation}
    \end{center}
    \caption{Ancilla augmentation of an elementary \ComputationActivationNetwork{} a) by replacing each leg vector of the elementary activation tensor by an directed ancilla augmentation.}
    \label{fig:canElementaryAugmentation}
\end{figure}

%% Elementary CAN Augmentation
Elementary \ComputationActivationNetworks{} can be augmented by ancilla augmentation of each leg vector in the activation tensor network (see \figref{fig:canElementaryAugmentation}).
We understand ancilla variables as variables in a Bayesian Network having single parents, corresponding with the ancilla augmentation of an elementary tensor.
For an example, see \figref{fig:toyAccounting}.

%% Hidden variables in activation graph
To a generic activation hypergraph $\graph$, we would do ancilla augmentation for the activation tensor network, where each hidden variable is treated as a distributed qubit.
To treat it as a distributed qubit, each hidden activation qubit is prepared in uniform state (i.e. a Hadamard gate acting on a ground state) and will be omitted in the measurement scheme (either not measured, or its measurement ignored).
This can be understood from a tensor network perspective, where marginalization means contracting, which is exactly how hidden activation variables are used.
For an example in the $\ttformat$ format, see \figref{fig:canTTAugmentation}.

\begin{figure}[t]
    \begin{center}
        \input{tikz_pics/sampling/can_TT_augmentation}
    \end{center}
    \caption{Ancilla augmentation of an $\ttformat$ \ComputationActivationNetwork{} a) by a Bayesian Network b).
    The hidden variables in the $\ttformat$ activation tensor are treated as distributed variables.
    Marginalization over the hidden variables and conditioning on the activation variables in state $1$, we reproduce the \ComputationActivationNetwork{}.}
    \label{fig:canTTAugmentation}
\end{figure}

\subsection{Amplitude Amplification}

The above scheme prepares the ancilla augmented \ComputationActivationNetwork{}
\begin{align*}
    \secprobat{\avariables,\headvariables,\shortcatvariables}
    = \acceptanceprob \cdot \onehotmapofat{1_{[\seldim]}}{\avariables}
    \otimes \contractionof{\bencodingofat{\sstat}{\headvariables,\shortcatvariables},\probwith}{\headvariables,\shortcatvariables}
    + (1-\acceptanceprob) \probofat{\perp}{\shortcatvariables,\headvariables,\avariables}
\end{align*}
and the prepared state is the combination of q-samples
\begin{align*}
    \qstateofat{0}{\avariables,\headvariables,\shortcatvariables}
    = \sqrt{\acceptanceprob} \cdot \onehotmapofat{1_{[\seldim]}}{\avariables}
    \otimes \contractionof{\bencodingofat{\sstat}{\headvariables,\shortcatvariables},\sqrt{\probwith}}{\headvariables,\shortcatvariables}
    + \sqrt{(1-\acceptanceprob)} \sqrt{\probofat{\perp}{\shortcatvariables,\headvariables,\avariables}} \, .
\end{align*}

Sampling from the augmented \ComputationActivationNetwork{} can be also done classically by forward sampling:
Just draw the variables from uniform, calculate the value qubit by a logical circuit inference and accept with probability by the computed value.

To achieve a quantum advantage we need amplitude amplification on the ancilla qubits.
%This can provide a square root speedup in the complexity compared with classical rejection sampling.

In the subspace $\subspaceof{}$ spanned by the axes
\begin{itemize}
    \item $\qstateof{\goodstatesymbol}=\onehotmapofat{1_{[\seldim]}}{\avariables}\otimes \contractionof{\bencodingofat{\sstat}{\headvariables,\shortcatvariables},\sqrt{\probwith}}{\headvariables,\shortcatvariables}$
    \item $\qstateof{\badstatesymbol}=\sqrt{\probof{\perp}}$
\end{itemize}
which can be thought of the polar angle coordinate of the state
\begin{align*}
    \qstateofat{0}{\avariables,\headvariables,\shortcatvariables}
    = \sinof{\frac{\rotanglesymbol}{2}} \qstateofat{0}{\avariables,\headvariables,\shortcatvariables}
    + \cosof{\frac{\rotanglesymbol}{2}} \qstateofat{1}{\avariables,\headvariables,\shortcatvariables}
\end{align*}
where we define an angle $\rotanglesymbol$ by
\begin{align*}
    \sinof{\frac{\rotanglesymbol}{2}} \coloneqq \sqrt{\acceptanceprob}
\end{align*}

The amplitude amplification scheme now iterates between
\begin{itemize}
    \item Reflection $\reflectionof{\goodstatesymbol}$ along $\qstateof{\goodstatesymbol}$, by
    \begin{align*}
        \reflectionofat{\goodstatesymbol}{\cavariableof{[\seldim]}{\insymbol},\cavariableof{[\seldim]}{\outsymbol}}
        = \identityat{\cavariableof{[\seldim]}{\insymbol},\cavariableof{[\seldim]}{\outsymbol}}
        - 2 \onehotmapofat{\onetuple{[\seldim]}}{\cavariableof{[\seldim]}{\insymbol}}
        \otimes \onehotmapofat{\onetuple{[\seldim]}}{\cavariableof{[\seldim]}{\outsymbol}}
    \end{align*}
    \item Reflection $\reflectionof{}$ along $\qstate$, by
    \begin{align*}
        \reflectionofat{\qstate}{\cavariableof{[\seldim]}{\insymbol},\cheadvariableof{[\seldim]}{\insymbol},\cavariableof{[\seldim]}{\insymbol},\cavariableof{[\seldim]}{\outsymbol},\cheadvariableof{[\seldim]}{\outsymbol},\cavariableof{[\seldim]}{\outsymbol}}
        =& \identityat{\cavariableof{[\seldim]}{\insymbol},\cheadvariableof{[\seldim]}{\insymbol},\cavariableof{[\seldim]}{\insymbol},\cavariableof{[\seldim]}{\outsymbol},\cheadvariableof{[\seldim]}{\outsymbol},\cavariableof{[\seldim]}{\outsymbol}} \\
        &- 2 \qstateofat{0}{\cavariableof{[\seldim]}{\insymbol},\cheadvariableof{[\seldim]}{\insymbol},\ccatvariableof{[\catorder]}{\insymbol}} \otimes
        \qstateofat{0}{\cavariableof{[\seldim]}{\outsymbol},\cheadvariableof{[\seldim]}{\outsymbol},\ccatvariableof{[\catorder]}{\outsymbol}}
        %U (\identityat{\cavariableof{[\seldim]}{\insymbol},\cavariableof{[\seldim]}{\outsymbol}} - 2 \onehotmapofat{1_{[\seldim]}}{\cavariableof{[\seldim]}{\insymbol}} \otimes \onehotmapofat{1_{[\seldim]}}{\cavariableof{[\seldim]}{\outsymbol}}) U^T
    \end{align*}
    This reflection can be done by concatenating the operator preparing the ancilla-augmentation of the \ComputationActivationNetwork{} with a initial state reflection.
\end{itemize}

Now the application of $\repnum$ grover operators $\reflectionof{}\reflectionof{\goodstatesymbol}$ on $\qstateof{0}$ prepares a state (see \figref{fig:rotationSketch})
\begin{align*}
    \qstateofat{\repnum}{\avariables,\headvariables,\shortcatvariables}
    = \sinof{\left(\frac{1}{2}+\repnum\right)\rotanglesymbol} \qstateofat{\goodstatesymbol}{\avariables,\headvariables,\shortcatvariables}
    + \sinof{\left(\frac{1}{2}+\repnum\right)\rotanglesymbol} \qstateofat{\badstatesymbol}{\avariables,\headvariables,\shortcatvariables}
\end{align*}
The state is closest to $\qstateof{\goodstatesymbol}$ when
\begin{align*}
    \left(\frac{1}{2}+\repnum\right)
    \rotanglesymbol \approx \frac{\pi}{2} \, .
\end{align*}

Estimating for small $\acceptanceprob$ (tight for small $\acceptanceprob$)
\begin{align*}
    \sqrt{\acceptanceprob} = \sinof{\frac{\rotanglesymbol}{2}}\leq \frac{\rotanglesymbol}{2}
\end{align*}
we get a bound
\begin{align*}
    \repnumof{\mathrm{opt}} \leq \frac{\pi}{4\sqrt{\acceptanceprob}} \, .
\end{align*}
This shows that we have a square root improvement on classical rejection sampling, which would require $\frac{1}{\acceptanceprob}$ repetitions to produce a sample.
Note however that we compare quantum circuit length with the number of classical repetitions in sampling.

\begin{figure}[t]
    \begin{center}
        \input{tikz_pics/sampling/rotation_sketch.tikz}
    \end{center}
    \caption{Sketch of the amplitude amplification scheme.}
    \label{fig:rotationSketch}
\end{figure}


%Note, that the variable qubits are uniformly distributed when only the computation circuit is applied.
%When sampling the probability distribution, we need the ancilla qubits to be in state $1$ in order for the sample to be valid.
%Any other states will have to be rejected.

%\textbf{Open Question:} Is there a way to avoid amplitude amplification and use a more direct circuit implementation of the activation network?
%- Cannot be the case, when the encoding is determined by the activation tensor alone: Needs to use the computated statistic as well.
%Negative result in \cite[Section~6.6]{nielsen_quantum_2010}: When using multiple applications of \computationCircuit{} in combination with further unitaries, need at least as many applications as with amplitude amplification.

\subsection{Sampling from \ComputationActivationNetworks{}}

For \ComputationActivationNetworks{} we find a q-sample by concatenating \computationCircuits{} and \activationCircuits{} (see \figref{fig:caCircuit}).

\begin{figure}
    \begin{center}
        \input{tikz_pics/sampling/ca_circuit}
    \end{center}
    \caption{
        Quantum Circuit to reproduce a q-sample of an augmented \ComputationActivationNetwork{}.
        We measure the distributed qubits $\shortcatvariables$ and the ancilla qubits $\avariableof{[\seldim]}$ and reject all samples, where an ancilla qubit is measured as $0$.
    }\label{fig:caCircuit}
\end{figure}

\input{examples/sampling/toy_accounting_circuit}

